% Problem record op_644f1aced78f36c9. This TeX file is derived from
% database/problems_json/op_644f1aced78f36c9.json by scripts/sync-tex.mjs; edit the JSON record
% and rerun the script rather than editing this file.
\section{Two-error perfect quantum codes over non-prime-power local dimensions}
\label{sec:op_644f1aced78f36c9}

\paragraph{Problem.}
Does a nontrivial pure perfect quantum code correcting arbitrary errors on two
qudits exist for some local dimension $q$ with at least two distinct prime
factors?  Let $q\geq2$ and $n\geq1$ be integers, and let
$\mathcal C\subseteq(\mathbb C^q)^{\otimes n}$ be a code of integer dimension
$K>1$ with orthogonal projector $P$.  No stabilizer structure is assumed, and
$K$ need not be a power of $q$.  On one site let
$X\lvert j\rangle=\lvert j+1 \bmod q\rangle$ and
$Z\lvert j\rangle=\omega^j\lvert j\rangle$ with $\omega=e^{2\pi i/q}$, and use
the phase-free Weyl basis $\{X^aZ^b: 0\leq a,b\leq q-1\}$ of $q^2$ operators,
including the identity.  Let $\mathcal E_2$ be the set of $n$-fold tensor
products of these operators with weight at most two, where weight counts
non-identity sites.  The code is a pure two-error-correcting code when
\begin{equation}
  PE^\dagger FP=\delta_{E,F}\,P,
  \qquad E,F\in\mathcal E_2,
  \label{eq:644f-pure}
\end{equation}
where $\delta_{E,F}$ is the Kronecker delta, so that distinct correctable
errors map the code space to mutually orthogonal subspaces.  Put
\begin{equation}
  V_2(n,q):=\lvert\mathcal E_2\rvert
  =1+n(q^2-1)+\binom n2(q^2-1)^2 .
  \label{eq:644f-volume}
\end{equation}
Equation~\eqref{eq:644f-pure} implies $K\,V_2(n,q)\leq q^n$, and the code is
Hamming-perfect when
\begin{equation}
  K\,V_2(n,q)=q^n .
  \label{eq:644f-perfect}
\end{equation}
The question asks whether Eqs.~\eqref{eq:644f-pure} and~\eqref{eq:644f-perfect}
hold simultaneously for some $n$, some integer $K>1$, and some $q$ that is not
a prime power.  A complete answer either exhibits such a code or proves that
none exists for any non-prime-power $q$.

\subsection*{Status}

Unsolved

\subsection*{Source}

Cazorla Garc\'{\i}a proposes, in Section~5.2 of the arXiv version, adapting his
Diophantine method for classical perfect two-error-correcting codes over
non-prime-power alphabets to perfect two-error-correcting quantum codes, noting
that the classification of Li and Xing assumes prime-power local dimension
\sourcecite{ref:644f-cazorla-garcia}{Caz24},
\sourcecite{ref:644f-li-xing}{LX13}.  The question is derived from this
documented gap; neither source states it as a named problem or conjecture.

\subsection*{Progress}

\begin{itemize}
  \item Every Weyl operator of weight between $1$ and $4$ is proportional to
  $E^\dagger F$ for some distinct $E,F\in\mathcal E_2$, so
  Eq.~\eqref{eq:644f-pure} forces distance at least five.  Rains's quantum
  Singleton bound $K\leq q^{n-2d+2}$ holds for every integer alphabet size
  (Theorem~2 of the arXiv version) and gives $K\leq q^{n-8}$.  Combined with
  Eq.~\eqref{eq:644f-perfect} and $K\in\mathbb Z$, it yields the necessary
  conditions
  \begin{equation}
    V_2(n,q)\geq q^8,\qquad V_2(n,q)\mid q^n,\qquad n\geq 9 .
    \label{eq:644f-necessary}
  \end{equation}
  These are direct deductions, not separate existence theorems
  \sourcecite{ref:644f-rains}{Rai99}.

  \item For prime-power $q$, Li and Xing prove a quantum analogue of Lloyd's theorem
  (Theorem~7) and show that the dimension of a pure perfect code is a power of
  $q$ (Lemma~8).  Theorem~9 of the arXiv version concludes that every
  nontrivial pure perfect quantum code has parameters
  $\bigl(\bigl((q^{2l}-1)/(q^2-1),\,q^{n-2l},\,3\bigr)\bigr)_q$, so minimum
  distance three, whether or not it is a stabilizer code.  Hence no pure perfect
  two-error-correcting code exists at prime-power local dimension.  This
  classifies parameters, not codes up to equivalence, and the prime-power
  hypothesis enters the proof.  The arXiv version omits the proof of Theorem~9
  as long, stating that it follows the classical Tiet\"av\"ainen--van Lint
  argument \sourcecite{ref:644f-li-xing}{LX13}.

  \item A hypothetical code yields a purely arithmetic consequence.  With
  $B_2(n,Q):=1+n(Q-1)+\binom n2(Q-1)^2$, the integers $Q:=q^2$ and $M:=Kq^n$
  satisfy $B_2(n,Q)=V_2(n,q)$ and hence, by Eq.~\eqref{eq:644f-perfect},
  \begin{equation}
    M\,B_2(n,Q)=Q^n ,
    \label{eq:644f-transfer}
  \end{equation}
  the classical sphere-packing equation for a perfect two-error-correcting code
  of $M$ words of length $n$ over $Q$ symbols.  This deduction does not produce a
  classical code, so a classical nonexistence theorem transfers automatically
  only when its proof uses Eq.~\eqref{eq:644f-transfer} alone.  Cazorla Garc\'{\i}a studies
  exactly this Diophantine equation in Sections~2 and~3
  \sourcecite{ref:644f-cazorla-garcia}{Caz24}.

  \item Lemma~3.1 and Table~2 of the arXiv version of that paper list all solutions
  of Eq.~\eqref{eq:644f-transfer} with $n\geq5$ for every non-prime-power
  $Q\geq6$ with either $Q\leq200$ and $Q\notin\{94,166\}$, or $Q\leq600$ and all
  prime factors of $Q$ in $\{2,3,5,7,11\}$; solutions occur only for
  $Q\in\{15,21,46\}$.  For every non-prime-power $q\leq24$ the square $Q=q^2$
  lies in this range and is not in that set, and $n\geq9$ by
  Eq.~\eqref{eq:644f-necessary}.  The transfer therefore excludes pure perfect
  two-error-correcting quantum codes for
  \begin{equation}
    q\in\{6,10,12,14,15,18,20,21,22,24\},
    \label{eq:644f-excluded}
  \end{equation}
  which are all non-prime-power integers up to $24$.  This list is a deduction
  from the cited lemma, not a quantum theorem stated in the source.  It uses only
  the lemma, not the paper's classical nonexistence theorem, whose proof also
  invokes Lloyd's theorem \sourcecite{ref:644f-cazorla-garcia}{Caz24}.

  \item Bennett proves that no classical perfect two-error-correcting code exists
  over a non-prime-power alphabet whose largest prime factor is at most $13$
  (Theorem~4).  For alphabet size $2^\alpha p^\beta$ with $p$ an odd prime and
  integers $\alpha,\beta\geq1$, such a code requires $\alpha>20$,
  $\beta\leq2519$, $p>10^{10}$, and $p\equiv3\pmod 8$ (Theorem~2).  Both
  proofs combine the sphere-packing condition with integer roots of the
  classical Lloyd polynomial, so they do not transfer to quantum codes through
  Eq.~\eqref{eq:644f-transfer} alone, and the preprint does not treat quantum
  codes \sourcecite{ref:644f-bennett}{Ben26}.  A subsequent preprint
  characterizes classical perfect codes as exact minimizers of a quadratic
  discrepancy and records that existence of perfect two-error-correcting codes
  over non-prime-power alphabets remains open; the characterization is not an
  existence proof \sourcecite{ref:644f-zabokritskiy}{Zab26}.

  \item Historical GitHub report (2026-09-24): Yuxuan Zhang (GitHub: yuxuanzhang1995) withdrew an earlier partial-result submission and asked that it not be treated as a current result; this link records the withdrawal only. \href{https://github.com/Naixu-Guo/quantum-open-problems/issues/107}{Issue \#107}.
\end{itemize}

\subsection*{References}

\begin{enumerate}[label={},leftmargin=4.8em,labelsep=0.5em]
  \footnotesize
  \item[\textup{[Rai99]}]\label{ref:644f-rains}
  E. M. Rains, ``Nonbinary quantum codes,'' \emph{IEEE Transactions on
  Information Theory} \textbf{45}(6), 1827--1832 (1999).
  \href{https://doi.org/10.1109/18.782103}{doi:10.1109/18.782103};
  \href{https://arxiv.org/abs/quant-ph/9703048}{arXiv:quant-ph/9703048}.

  \item[\textup{[LX13]}]\label{ref:644f-li-xing}
  Z. Li and L. Xing, ``Classification of $q$-ary perfect quantum codes,''
  \emph{IEEE Transactions on Information Theory} \textbf{59}(1), 631--634 (2013).
  \href{https://doi.org/10.1109/TIT.2012.2217475}{doi:10.1109/TIT.2012.2217475};
  \href{https://arxiv.org/abs/0907.0049}{arXiv:0907.0049}.

  \item[\textup{[Caz24]}]\label{ref:644f-cazorla-garcia}
  P.-J. Cazorla Garc\'{\i}a, ``Perfect codes over non-prime power alphabets: an
  approach based on Diophantine equations,'' \emph{Mathematics} \textbf{12}(11),
  1642 (2024).
  \href{https://doi.org/10.3390/math12111642}{doi:10.3390/math12111642};
  \href{https://arxiv.org/abs/2405.03347}{arXiv:2405.03347}.

  \item[\textup{[Ben26]}]\label{ref:644f-bennett}
  M. A. Bennett, ``Perfect $2$-codes over arbitrary alphabets,'' arXiv preprint
  (2026).
  \href{https://arxiv.org/abs/2607.27555}{arXiv:2607.27555}.

  \item[\textup{[Zab26]}]\label{ref:644f-zabokritskiy}
  A. L. Zabokritskiy (Yohananov), ``Perfect codes as exact minimizers of
  quadratic discrepancy in $q$-ary Hamming spaces,'' arXiv preprint (2026).
  \href{https://arxiv.org/abs/2608.01134}{arXiv:2608.01134}.
\end{enumerate}

\subsection*{Comment}

No construction and no nonexistence proof covering all non-prime-power local
dimensions was located.  The exclusions above cover only the dimensions in
Eq.~\eqref{eq:644f-excluded}, so the smallest non-prime-power dimension they
leave undecided is $q=26$.  The prime-power classification, the arithmetic
exclusions, and the classical results above do not classify the quantum codes
considered here.  Bennett's classical arguments use Lloyd's theorem at alphabet
size $Q$.  The quantum analogue of Li and Xing, whose Lloyd polynomial coincides
with the classical one at $Q=q^2$, is proved under the prime-power hypothesis,
so extending those exclusions to non-prime-power $q$ requires that ingredient.
Arithmetic feasibility of Eqs.~\eqref{eq:644f-necessary}
and~\eqref{eq:644f-transfer} would not by itself construct a code.  Literature
checked through 15 September 2026.

\subsection*{Field}

Quantum Error Correction

\subsection*{Topic}

Quantum coding theory

\subsection*{ID}

\texttt{op\_644f1aced78f36c9}
